\documentclass[11pt]{article}

%------------------------------------------------------------------------
\usepackage{geometry}
\usepackage{latexsym}
\usepackage[empty]{fullpage}
\usepackage{titlesec}
\usepackage{parskip}
\usepackage{marvosym}
\usepackage[usenames,dvipsnames]{color}
\usepackage{verbatim}
\usepackage{enumitem}
\usepackage[hidelinks]{hyperref}
\usepackage{fancyhdr}
\usepackage{lineno}
\usepackage{lastpage}
%\usepackage{lipsum}
\usepackage{color}
\usepackage{setspace}
\usepackage{textcomp}
\usepackage{ragged2e}
%\usepackage{hyperref}
%\usepackage[bitstream-charter]{mathdesign} % for journal-like fonts
\usepackage{url}
\hypersetup{colorlinks,linkcolor=blue,urlcolor=blue}
\urlstyle{ttfamily}
\usepackage[superscript]{cite}
\usepackage{microtype}

\renewcommand\citeform[1]{[#1]}
\newcommand\rt[2] {{\textcolor{red}{#1}}{\textcolor{green}{#2}}}

%------------------------------------------------------------------------
% Paper margin setup
 \geometry{
 letterpaper,
 total={216mm,280mm},
 left=25mm,
 right=25mm,
 top=22mm,
 bottom=25mm
 }
 
%-----------------------------------------------------------------------
\pagestyle{fancy}
\fancyhf{} 

\lfoot[<even output>]{\color{Gray}Teaching philosophy}
\cfoot[<even output>]{\color{Gray}Page \thepage\ of \pageref{LastPage}}
\rfoot[<even output>]{\color{Gray}S. Sunoj}

\renewcommand{\headrulewidth}{0pt}
\renewcommand{\footrulewidth}{0pt}

\urlstyle{same}

%\raggedbottom
%\raggedright
%\setlength{\tabcolsep}{1in}
%------------------------------------------------------------------------
\newcommand\idea[1]{{\bfseries #1}\\}

\titleformat{\section}
%  {\normalfont\bfseries}{\thesection.}{0.5em}{\vspace{-1.5em}} 
  {\normalfont\bfseries}{\thesection.}{0.4em}{\vspace{-0.8em}} 
  
\setlength{\parindent}{2em}

\begin{document}
\setstretch{1.1}

\begin{center}
\hrule
%\vspace{0.5em}
{\Large\textbf{Statement of teaching philosophy}}\\
\vspace{0.5em}
\normalsize{Sunoj Shajahan\\
Postdoctoral Associate, Nutrient Management Spear Program,\\
Cornell University, Ithaca, NY.}
\vspace{0.5em}
\hrule
\end{center}
\vspace{0.5em}

\begin{center}
\textls[-20]{\emph{``A good teacher can inspire hope, ignite the imagination, and instill a love of learning'' --- Brad Henry}}
\end{center}

\justify
{\hspace{1em} The teaching profession makes every other profession possible. My passion for teaching started during my school days when my science teacher assigned me to tutor a few students and since then I have developed an ardent passion for teaching. I always strive to give the best efforts in teaching, irrespective of the size of the group. I was fortunate to have extraordinary teachers throughout my education, who have been my inspiration. Additionally, I have taken part in a few teaching workshops at Cornell University and equipped myself with effective teaching strategies. My teaching philosophy includes the following five strategies:

\begin{enumerate}[itemsep = -1pt]
\item Providing differential instructions for students with different learning potential 
\item Using analogy to teach complex concepts
\item Preparing the lecture materials effectively by emphasizing keywords and images
\item Inviting students to think critically
\item Encouraging students to build effective communication skills (oral and written)
\end{enumerate}

%\vspace{-0.5em}
%\section{Providing differential instructions for students with different learning potential }
%\hspace{1em} I understand any classroom will be a diverse environment. Different students have different learning potential and understanding the student's potential is one of the significant aspects of teaching. This links to the student's prior knowledge on the subject. Therefore, at the beginning of the course, I will conduct a preliminary quiz to evaluate the student's knowledge on the subject. I will follow active learning strategies through group presentations, group activities, discussions, debates, to encourage more class interactions from the students. 
%
%\vspace{-0.5em}
%\section{Using analogy to teach complex concepts}
%\hspace{1em} I encountered an interesting experience during my participation at the 3 Minute Thesis (3MT) competition during my Ph.D. I was initially mind-stuck during the preparation for the competition since I had to explain my research (based on image processing algorithms) to the audiences from unfamiliar research fields with a time constraint of three minutes but soon realized the key was to \textit{explain my research in the form of a simple story} that everyone can relate to. Explaining the concepts with day-to-day examples made it easier. Therefore, I have decided to make this as one of my key teaching practices and I still apply this to most of my presentation talks. 
%
%%Also, technical jargons might be quite intimidating for students. In precision agriculture, new and upcoming terminologies in big data, internet of things, and machine learning sound technical; however, this can be explained in a simple non-technical language. 
%
%%\hspace{1em} For example, machine learning idea can be related to how a toddler is trained in colors and shapes. We repeatedly display and say the names of the colors and shapes to the toddler until he/she is trained. Here the `toddler’ can be compared to the `machine’, `colors and shapes’ to the `feature’, `teaching toddler’ to `training’ the machine, and `repeated display and saying names’ to the actual `annotation’ process in machine learning. This is the overarching idea of machine learning. Explaining concepts to students with such examples makes it easier for the students not only to understand but also to remember for a long time. Hence, I will work on simplifying newer and complex concepts by following this practice.
%
%\vspace{-0.5em}
%\section{Preparing the lecture materials effectively by emphasizing keywords and images}
%\hspace{1em} I believe that teaching is much more than just delivering knowledge in a convenient format to the students. Ideally, there is an expectation from the students to have a few take-home points from the lecture. To fulfill this, \emph{I aim to effectively prepare lecture materials using visual aids and emphasizing keywords and images}. I strive to prepare the slides with minimal text, more images, and highlighting keywords. At the end of each lecture, \emph{key take-home points} will be discussed with the students. Also, before the start of a class, I would spend a few minutes reviewing the concepts dealt with in the previous class session. With my earlier teaching experience, I found this practice to be more effective and efficient to better engage the students.
%
%\newpage
%\vspace{-0.5em}
%\section{Inviting students to think critically}
%\hspace{1em} Having come from an agricultural engineering background I developed a fascination with computer programming because it strengthened logical thinking and mathematical concepts. Developing logical algorithms and fixing coding errors helped me to critically think and improve my problem-solving capabilities. I aim to teach and train the students to program and develop their own software tools and modeling programs using open-source software platforms. I remember the following quote by Steve Jobs – ``\emph{Everyone should know how to program a computer because it teaches you how to think}”. I am confident that through this approach, I will provide the students a sense of confidence to critically think. 
%%Making students to critically think aligns with Iowa State University's President Wintersteen's focus \href{https://www.diversity.iastate.edu/}{(Link)}. 
%
%\vspace{-0.5em}
%\section{Encouraging students to build effective communication skills (oral and written)}
%\hspace{1em} I will let them choose their own ideas for the term paper related to the course and ask them to provide biweekly progress through presentations, posters, or a 3-minute talk, which provides a platform for the students to prepare and present. This classroom exercise will improve their creativity and confidence in effectively communicating ideas. For writing, I will train them to effectively communicate through writing by assigning individual project reports, group reports, etc. I aim to offer timely feedback and suggestions when necessary by conducting one-to-one meetings with the students. This way, I can have face-to-face interactions with the students, which will hopefully make them feel more comfortable approaching me even after the scheduled class hours. 
%
%I will encourage students to participate in technical conferences and college-level symposiums to present their project/ideas by sharing them with the details of such events, bonus points will be assigned for such activities. 
%I will make it a point to arrive at least ten minutes early to class as it provides an opportunity for the students to discuss assignments, doubts, and their needs if any. 

%\vspace{-1em}
\hspace{1em} I have followed some of these strategies during my previous experience as a teaching instructor in Ramakrishna Mission Vivekananda Educational and Research Institute, Coimbatore, India. I handled an undergraduate-level course ``Principles of food science and processing for BS Agriculture'' (class strength - 60 students). During the course, \emph{I took efforts to remember all the students' names} as I believe that addressing students by their correct name establishes a good connection between the student and the instructor. I also believe that involving the students in class discussion activities is one of the essentials in learning (“\emph{Tell me and I forget, teach me and I remember, involve me and I learn}”~---~Benjamin Franklin). 

%Furthermore, I collected biweekly feedback from student, which significantly helped me to improve my teaching and to solve some of the issues faced by the students; hence, I will practice these in future too. Finally, I will keep track and account for the progress made by each student as it should be recognized. I made explicit attempts to follow the same in the past. All these efforts have made me an approachable instructor and the students enjoyed attending my class.

\vspace{1em}
\section*{Teaching interests at University of Illinois at Urbana-Champaign}
\vspace{-0.8em}
\hspace{1em} 
My educational background in food engineering and remote sensing, along with interdisciplinary research training has prepared me to teach and handle a wide variety of courses at the undergraduate and graduate level in the Department of Agricultural and Biological Engineering at the University of Illinois at Urbana-Champaign including AGRON 202 (Intro to Precision Ag Software), AGRON 655 (Site-Specific Agriculture). Further, I will be interested in developing new courses at the undergraduate and graduate levels such as \emph{Geospatial Technologies in Agriculture and Environment}, \emph{Introduction to Google Earth Engine}, \emph{Introduction to data analysis in agriculture}, and \emph{Applications of big data in agriculture}. Also, my philosophy of using free and open-source software (FOSS) will allow students to gain and apply their knowledge even during unprecedented situations, like the COVID-19 pandemic. 

%I would be able to develop and handle new courses, which would involve digital image processing and the basics of programming for agriculture-related majors. 

%I am enthusiastic about supporting students through computer programming and data analytics in the Department of Agricultural and Biosystems Engineering at NDSU.
%
%
%or other courses depending on departmental needs. Additionally, with my experience, I will be able to teach or co-teach other courses offered in ABEN such as ABEN 255. Computer Aided Analysis \& Design, ABEN 263. Biological Materials Processing, ABEN 377. Numerical Modeling in Agricultural and Biosystems Engineering, ABEN 444. Transport Processes, ABEN 458. Process Engineering for Food, Biofuels and Bioproducts, ASM 454. Principles and Application of Precision Agriculture, and ASM 455. Data Management in Precision Agriculture. Based on my diverse experiences, skills, and interests, I am enthusiastic about supporting students through computer programming and data analytics in the Department of Agricultural and Biosystems Engineering at NDSU.


%I have my bachelors and masters in Agricultural Processing and Food Engineering and a Ph.D. in Agricultural and Biosystems Engineering with focus on precision agriculture. With this mix of experiences, I will be able to handle the following courses: ABE 160: Systematic Problem Solving and Computer Programming, ABE 216: Fundamentals of Agricultural and Biosystems Engineering, ABE 451: Food and Bioprocess Engineering, ABE 506: Applied Computational Intelligence, ABE 556: GIS Programming and Automation, ABE 569: Engineering for Grain Storage, Preservation, Handling, and Processing Systems. 

%I will be interested in developing a ``Precision agriculture data analysis'' course within ABE that will have more of hands-on experience with exposure to different data sources, analysis such data, and developing models or tools for those datasets. I will also make sure it is different from the TSM 433 and TSM 533 offered in Technology Systems Management.      


\end{document}